\documentclass[preprint,12pt,authoryear]{elsarticle}

\usepackage{amsmath,amssymb,amsthm}
\usepackage{algorithm}
\usepackage{algpseudocode}
\usepackage{booktabs}
\usepackage{multirow}
\usepackage{graphicx}
\usepackage{subcaption}
\usepackage{microtype}
\usepackage{xcolor}
\usepackage{float}          % [H] placement
\usepackage{placeins}       % \FloatBarrier between sections
\usepackage[hidelinks]{hyperref}
\usepackage{tikz}
\usetikzlibrary{arrows.meta, shapes.geometric, positioning, fit, backgrounds}

\pagestyle{empty}           % suppress page numbers in draft

%% Theorem environments
\newtheorem{proposition}{Proposition}
\newtheorem{corollary}{Corollary}
\newtheorem{remark}{Remark}

%% Notation shortcuts
\newcommand{\J}{\mathcal{J}}           % facilities
\newcommand{\I}{\mathcal{I}}           % clients
\newcommand{\K}{\mathcal{K}}           % kernel
\newcommand{\Jone}{J_1}               % LP-open
\newcommand{\Jfrac}{J_{\text{frac}}}  % LP-fractional
\newcommand{\Jzero}{J_0}              % LP-closed
\newcommand{\xLP}{x^{\text{LP}}}
\newcommand{\yLP}{y^{\text{LP}}}
\newcommand{\rc}[1]{\bar{c}_{#1}}     % reduced cost

\begin{document}

\begin{frontmatter}

\title{Revisiting Kernel Search for the Large-Scale Multi-Source
       Capacitated Facility Location Problem:
       Three Targeted Design Improvements}

\author[jmu]{Kamal Lamsal}
\ead{lamsalkx@jmu.edu}

\address[jmu]{Department of Computer Information Systems and Business Analytics,
              James Madison University, Harrisonburg, Virginia, USA}

\begin{highlights}
\item We revisit Kernel Search for the large-scale Multi-Source CFLP and
      separate its runtime into LP and restricted-MILP components.
\item Column generation targets the LP-relaxation bottleneck while preserving
      the LP-derived information used by KS.
\item A flat-$k$ edge rule gives explicit base edge-set sizes across instance
      families and capacity regimes.
\item Preserving incumbent arcs between stages makes the previous incumbent
      feasible in the next restricted MILP.
\item The computational study is structured to isolate LP acceleration,
      MILP-stage acceleration, and end-to-end performance.
\end{highlights}

\begin{abstract}
Kernel Search (KS) remains one of the strongest matheuristics for the
large-scale Multi-Source Capacitated Facility Location Problem (MSCFLP).
This paper takes a design-analysis perspective on G\&S~2012 Kernel Search.
A close examination of the original method reveals three design choices that
can each be refined independently, yielding consistent improvements in both
solution quality and computation time.

The first concerns LP construction cost: the full LP relaxation contains all
$|\J||\I|$ shipping variables, most of which play no role in the optimal
solution.
Column generation recovers the same bound, facility partition, and
reduced-cost signal at substantially lower cost.

The second concerns arc selection: the global reduced-cost threshold $\gamma$
used to build restricted MILPs provides no size guarantee and systematically
admits fewer arcs for bucket facilities than for kernel facilities, because
bucket facilities have LP dual zero and receive no capacity-cost discount on
their arc reduced costs.

The third concerns stage transitions: when the arc set changes between stages,
the incumbent solution from the previous subproblem may no longer be feasible
in the next, because arcs it uses may be absent.
The solver must rediscover feasibility from scratch at every stage transition.

We address each with a targeted modification.
Column generation replaces full LP construction.
A client-centric flat-$k$ criterion replaces the global $\gamma$ threshold,
selecting exactly $k$ arcs per client by reduced cost and treating all
facilities uniformly.
Stage transitions are made feasibility-preserving by including every arc used
by the current incumbent when building the next edge set; we prove that this
guarantees the incumbent is feasible in the next restricted MILP by
construction, making it a valid warm start at every stage.
The resulting repaired method retains the full architecture of G\&S~2012 and
is named \emph{Funnel Kernel Search} (FKS).

The computational study isolates each modification via a dedicated comparison,
so that each improvement can be attributed to a specific design decision.
Results are reported across five large-scale benchmark families from the
Avella et al.\ test bed.
FKS finds the proven optimal solution on 109 of the 110 Test Bed~A instances
for which an optimum is known --- solutions that the state-of-the-art Benders
decomposition of Fischetti et al.\ requires up to 50{,}000 seconds per instance
to certify --- while FKS operates within a 300-second per-subproblem budget.
On Test Bed~B and Test Bed~C, FKS establishes 22 new best-known solutions,
improving on all published upper bounds it encounters.

\end{abstract}

\begin{keyword}
Capacitated facility location \sep Kernel search \sep Column generation \sep
Matheuristic \sep LP relaxation \sep Feasibility preservation \sep Multi-source
\end{keyword}

\end{frontmatter}

\section{Introduction}
\label{sec:intro}

The \emph{Capacitated Facility Location Problem} (CFLP) is a classical model for
selecting a set of facilities and allocating customer demand to them under
capacity restrictions.
It appears in distribution network design, supply chain planning, emergency
service location, and telecommunications infrastructure
design~\citep{klose2005facility,daskin2011network,laporte2019location}.
In the \emph{Multi-Source} CFLP (MSCFLP), a customer's demand may be split
among several open facilities.
In the \emph{Single-Source} CFLP (SSCFLP), each customer must be assigned to one
facility; this additional assignment restriction substantially changes the
structure and computational difficulty of the problem~\citep{klose2005facility}.
Both variants are NP-hard~\citep{sridharan1995capacitated}.

This paper focuses on the large-scale MSCFLP.
Avella et al.~\citep{avella2008effective} established a widely used
computational benchmark for this setting and proposed a Lagrangean-relaxation
and core-problem heuristic for instances with millions of variables.
Guastaroba and Speranza~\citep{guastaroba2012kernel} introduced a
Kernel Search (KS) matheuristic for the same problem class
(which they refer to simply as the CFLP; we retain the MSCFLP label to
distinguish it from the single-source variant).
Their method solves the LP relaxation, identifies an initial kernel of
LP-open and LP-fractional facilities, and solves a sequence of restricted MILPs
that add promising LP-closed facilities through capacity-expanding buckets.
The approach is effective because the strong MSCFLP formulation includes
variable upper-bound constraints, which make the LP solution a reliable predictor
of the facilities and arcs used in high-quality integer solutions.
Speed controls in the original paper act on the restricted-MILP phase: bucket
count limits, variable-count caps per subproblem, a kernel-pruning rule based on
repeated non-selection, and a forcing constraint requiring each bucket subproblem
to use a facility from the current bucket.
A faster variant, B-KS(0), searches only over the initial LP-selected kernel.
The original paper frames KS performance as a balance between restricted-search
quality and the time required to solve the resulting MILPs.

On the exact-method side, Görtz and Klose~\citep{gortz2012lagrangian} proposed a
Lagrangian branch-and-bound that is competitive with the best exact approaches on
the smaller GK instances.
Fischetti, Ljubić, and Sinnl~\citep{fischetti2016benders} later showed that a
modern Benders decomposition embedded in a branch-and-cut solver proves optimality
for 210 out of 445 large-scale p* instances within a 50{,}000-second time limit,
and consistently produces smaller gaps than KS across all three test beds.
Their computational study confirms that solving the LP at the root node
dominates the total runtime for the largest instances: more than 90\% of overall
computing time may be spent in LP-based cut separation.
To the best of our knowledge, these are the only results available for exact
methods on the large-scale multi-source CFLP benchmark of Avella et al.

This paper takes a different starting point: rather than proposing a new
acceleration strategy, we ask whether the design choices embedded in G\&S~2012
are each well-founded.
A close reading of the method reveals three design choices, each of which can
be refined independently.

\textbf{LP construction cost.}
The full LP relaxation contains $|\J||\I|$ shipping variables.
At optimality, most carry positive reduced cost and contribute nothing to the
LP solution or to the KS guidance it produces.
Constructing the full matrix is therefore unnecessary: column generation recovers
the same bound, facility partition, and reduced-cost signal without ever
materialising the irrelevant variables.

\textbf{Arc selection bias.}
G\&S~2012 includes arc $(j,i)$ whenever its reduced cost falls below a global
threshold $\gamma$, the mean reduced cost over LP-active arcs.
The threshold provides no guarantee on subproblem size and is structurally
biased: a bucket facility $j \in \Jzero$ has LP dual $\alpha^*_j = 0$, so its
arc reduced cost $c_{ji} - \pi^*_i$ contains no capacity signal.
A kernel facility with positive dual receives a capacity-dual discount that
lowers its reduced costs; the same $\gamma$ therefore admits more kernel arcs
than bucket arcs, giving bucket facilities a systematically smaller arc set
in every restricted MILP.

\textbf{Cold-start stage transitions.}
G\&S~2012 passes no solution information between consecutive restricted MILPs.
When the arc set is modified between stages, the incumbent from the previous
subproblem may no longer be feasible: arcs it uses may simply be absent from the
new edge set.
Because feasibility cannot be guaranteed, the incumbent cannot be supplied as
a warm start, and the solver must rediscover a feasible solution from scratch at
every stage.
This is not an obvious limitation of warm starting in general; it is a
structural consequence of arc-set transitions that discard incumbent arcs.

We address each in turn.
The LP is solved by column generation.
Arc selection is replaced by a client-centric flat-$k$ criterion: for each
client, the $k$ kernel facilities with lowest reduced cost are included,
giving an exact subproblem-size guarantee and treating all facilities uniformly
regardless of their LP dual value.
Stage transitions are made feasibility-preserving by including every arc used by
the current incumbent when building the next stage's edge set; we prove that
this construction guarantees the incumbent is feasible in the next restricted
MILP, making it a valid warm start by construction.
The resulting repaired variant is \emph{Funnel Kernel Search} (FKS).
It retains the full architecture of G\&S~2012 --- the strong MSCFLP formulation,
the LP-defined kernel, reduced-cost guidance, and bucket expansion --- and
modifies only the three components identified above.

The contributions are as follows.
\begin{itemize}
  \item \textbf{Design analysis.}
        We examine the G\&S~2012 design and identify three aspects where
        targeted modifications yield consistent improvements:
        (1)~the LP is built in full when column generation recovers the same
        information at lower cost;
        (2)~the global $\gamma$ threshold for arc selection admits no size
        guarantee and treats bucket facilities less favorably, since their LP
        dual is zero and their arc reduced costs carry no capacity signal,
        causing the same $\gamma$ to admit fewer arcs than for kernel
        facilities;
        (3)~the restricted-MILP phase starts cold at each stage, forcing the
        solver to rediscover feasibility after every arc-set transition.
  \item \textbf{Fix~1 --- LP cost.}
        We apply column generation to the MSCFLP LP relaxation,
        retaining the bound, facility partition, and reduced-cost signal
        required by KS while avoiding full-matrix construction.
  \item \textbf{Fix~3 --- feasibility-preserving stage transitions.}
        We prove that including every incumbent arc in the next stage's edge set
        guarantees the previous solution is feasible in the next restricted MILP
        (Proposition~\ref{prop:feasibility}).
        This is a structural result, not a heuristic: feasibility is guaranteed
        by construction and makes the incumbent a provably valid MIP start at
        every stage.
  \item \textbf{Fix~2 --- arc selection.}
        With the feasibility guarantee in place, we introduce a client-centric
        flat-$k$ criterion that replaces the global $\gamma$ threshold.
        The flat-$k$ rule narrows the fill set aggressively across stages
        (50$\to$20$\to$10 arcs per client) without risking cold restarts,
        because Proposition~\ref{prop:feasibility} holds regardless of how
        small $k$ becomes.
        It also gives an exact subproblem-size guarantee
        (Proposition~\ref{prop:size}) and treats all facilities uniformly.
  \item \textbf{Component-wise study.}
        The computational study isolates each fix via a dedicated comparison:
        full LP versus CG LP (Fix~1 alone), KS-CG versus KS-CG-WS (Fix~3
        alone), KS-CG-WS versus FKS-CG-WS (Fix~2 alone), and end-to-end
        KS-full versus FKS-CG-WS (combined effect).
\end{itemize}

The result is a matheuristic that is not only faster than the original but
also consistently finds solutions of higher quality. On the 110 Test Bed~A
instances with known proven optima, FKS finds the optimal solution on 109.

The remainder of the paper is organized as follows.
Section~\ref{sec:formulation} gives the MSCFLP formulation and the LP structure
used by KS.
Section~\ref{sec:algorithm} examines these three design aspects formally,
then presents the modifications in order of dependency: column-generation LP
(Fix~1), feasibility-preserving stage transitions with their structural proof
(Fix~3), and flat-$k$ arc selection (Fix~2) --- which relies on Fix~3 to
narrow arc sets safely.
Section~\ref{sec:experiments} describes the component-wise computational study
and reports results on the large-scale Test Bed benchmarks.
Section~\ref{sec:conclusion} concludes.

\FloatBarrier
\section{Problem Formulation and LP Structure}
\label{sec:formulation}

\subsection{The Multi-Source CFLP}

Let $\J = \{1,\ldots,m\}$ be the set of potential facility locations and
$\I = \{1,\ldots,n\}$ the set of clients.
Each facility $j \in \J$ has a fixed opening cost $f_j > 0$ and a capacity
$q_j > 0$.
Each client $i \in \I$ has a demand $d_i > 0$.
The unit shipping cost from facility $j$ to client $i$ is $c_{ji} \geq 0$.

Let $y_j \in \{0,1\}$ indicate whether facility $j$ is open, and let
$x_{ji} \in [0, d_i]$ be the amount of client $i$'s demand served by
facility $j$.
The MSCFLP is:

\begin{align}
  \min \quad & \sum_{j \in \J} f_j y_j
               + \sum_{j \in \J} \sum_{i \in \I} c_{ji} x_{ji}
               \label{eq:obj} \\
  \text{s.t.} \quad
  & \sum_{j \in \J} x_{ji} = d_i
    && \forall i \in \I
    \label{eq:demand} \\
  & \sum_{i \in \I} x_{ji} \leq q_j y_j
    && \forall j \in \J
    \label{eq:capacity} \\
  & x_{ji} \leq d_i y_j
    && \forall j \in \J,\; i \in \I
    \label{eq:vub} \\
  & y_j \in \{0,1\}
    && \forall j \in \J
    \label{eq:binary} \\
  & x_{ji} \geq 0
    && \forall j \in \J,\; i \in \I.
    \label{eq:nonneg}
\end{align}

Constraints~\eqref{eq:demand} ensure all demand is satisfied.
Constraints~\eqref{eq:capacity} enforce facility capacities.
The variable-upper-bound (VUB) constraints~\eqref{eq:vub} prevent routing
demand to closed facilities.
Together, constraints~\eqref{eq:capacity} and~\eqref{eq:vub} ensure that
$y_j = 0$ implies $x_{ji} = 0$ for all $i$, tightening the LP relaxation
substantially compared to formulations without VUB constraints.

The LP relaxation is obtained by replacing~\eqref{eq:binary} with
$0 \leq y_j \leq 1$.
We denote its optimal solution by $(\xLP, \yLP)$ and its objective value
by $z^{\text{LP}}$.
Unless explicitly stated otherwise, gaps reported in this paper are measured
relative to $z^{\text{LP}}$:
\[
  \text{gap}(\%) = \frac{z^H - z^{\text{LP}}}{z^{\text{LP}}} \times 100,
\]
where $z^H$ is the heuristic objective.
This gives a common lower-bound-based measure across all algorithmic variants.
When comparing with published tables, we also report objective values because
published gaps may use different lower bounds or hardware.

\subsection{LP Solution Structure}
\label{sec:lp_structure}

The LP solution partitions facilities into three disjoint sets based on
$\yLP_j$:
\begin{itemize}
  \item $\Jone = \{j \in \J : \yLP_j = 1\}$: LP-open facilities,
  \item $\Jfrac = \{j \in \J : 0 < \yLP_j < 1\}$: LP-fractional
        facilities,
  \item $\Jzero = \{j \in \J : \yLP_j = 0\}$: LP-closed facilities.
\end{itemize}
The \emph{kernel} is $\K = \Jone \cup \Jfrac$, following~\citet{guastaroba2012kernel}.

\paragraph{Why the kernel predicts integer solutions.}
The VUB constraints~\eqref{eq:vub} force $x_{ji} \leq d_i y_j$: no demand
can be routed to a facility without opening it, even fractionally, in the LP.
A LP-closed facility ($\yLP_j = 0$) therefore carries no LP flow; the LP
found it suboptimal to open facility $j$ even fractionally.
This is a useful signal: a facility that is LP-closed must overcome both fixed
opening cost and capacity interactions before it becomes attractive in an
integer solution.
The computational evidence in \citet{guastaroba2012kernel} shows that the
LP-defined kernel captures most facilities used in high-quality solutions on the
large-scale benchmark families.
This motivates the Kernel Search strategy: rather than solving the full
$m$-facility MILP, start from $\K$ and test promising LP-closed facilities
through restricted subproblems.

\paragraph{Reduced costs.}
Let $\pi_i \geq 0$ be the dual variable for demand constraint~\eqref{eq:demand}
and $\alpha_j \leq 0$ the dual variable for capacity
constraint~\eqref{eq:capacity}.
For the LP relaxation, the \emph{reduced cost} of shipping variable $x_{ji}$
is:
\begin{equation}
  \rc{ji} = c_{ji} - \pi_i - \alpha_j.
  \label{eq:rc}
\end{equation}
At LP optimality, $\rc{ji} \geq 0$ for every variable at its lower bound
($x_{ji} = 0$) and $\rc{ji} = 0$ for every basic variable ($x_{ji} > 0$).
A variable with large positive reduced cost is far from being optimal to
include: the LP would need a much cheaper shipping cost, or a much better
dual price, before it would use that arc.
This signal guides both column generation (Section~\ref{sec:cg}) and
edge selection (Section~\ref{sec:edge_selection}).

\paragraph{Why column generation is natural.}
The LP solution of the strong MSCFLP formulation is often sparse relative to the
complete $|\J||\I|$ shipping-variable matrix.
Many shipping variables are zero at LP optimality and have nonnegative reduced
cost.
Column generation exploits this structure by leaving such variables outside the
restricted master problem until their reduced cost indicates that they may
improve the LP objective.
Study~\ref{sec:study_lp} reports the observed active-column counts and checks
LP agreement against the full formulation where the full LP can be solved.

\FloatBarrier
\section{Design Analysis and Algorithm}
\label{sec:algorithm}

This section first recalls the G\&S~2012 baseline
(Section~\ref{sec:ks_baseline}), then examines the three design aspects
identified in the introduction (Section~\ref{sec:audit}).
The three modifications are presented in order of dependency:
column-generation LP (Section~\ref{sec:cg}),
feasibility-preserving stage transitions (Section~\ref{sec:proposition}),
and flat-$k$ arc selection (Section~\ref{sec:edge_selection}).
Fix~3 is presented before Fix~2 because Proposition~\ref{prop:feasibility}
is the structural reason the funnel's narrowing arc sets are safe to use:
the fill size $k$ can be reduced aggressively only because the incumbent
is always feasible by construction regardless of how small $k$ becomes.
Section~\ref{sec:full_algorithm} presents the complete FKS-CG-WS algorithm.

\subsection{Baseline Kernel Search (G\&S 2012)}
\label{sec:ks_baseline}

After solving the LP relaxation, facilities with positive LP dual
($\alpha^*_j > 0$) form the \emph{kernel} $\K = \Jone \cup \Jfrac$;
the remaining LP-closed facilities ($\alpha^*_j = 0$) form $\Jzero$ and are
organized into \emph{buckets}.
G\&S~2012 then proceeds as follows:

\begin{enumerate}
  \item \textbf{Kernel MILP.}
    Solve MILP restricted to facilities in kernel $\K$ and edges
    $\mathcal{E} = \{(j,i) : j \in \K,\; \rc{ji} \leq \gamma\}$,
    where $\gamma = \text{mean}\{\rc{ji} : \xLP_{ji} > 0\}$ is the mean
    reduced cost over LP-active arcs.
    An objective cutoff $z^H$ is imposed.
  \item \textbf{Bucket expansion.}
    Sort $\Jzero$ by $\rc{y_j}$ ascending.
    Add the next $l_{\text{buck}} = |\K|$ facilities as bucket $B_h$,
    connect to clients with $\rc{ji} \leq \gamma$, and solve
    MILP$(\K \cup B_h)$ with a forcing constraint $\sum_{j \in B_h} y_j \geq 1$.
    Repeat for up to $\bar{N}_B = 3$ buckets.
  \item \textbf{Kernel update.}
    If an improved solution opens facilities from the bucket, add them to the
    kernel.
    Facilities already in the kernel are removed after repeated non-selection,
    controlled by parameter $p$.
    An iterative variant restarts the solution phase when new facilities are
    added.
\end{enumerate}

This design gives several MILP-side controls over running time: the number of
buckets, their size, the removal rule, and the time limit per restricted MILP.

\subsection{Design Analysis: Three Refinement Opportunities}
\label{sec:audit}

Before presenting FKS, we examine three aspects of the G\&S~2012 design where
targeted modifications improve performance.
Each corresponds directly to one component of FKS.

\paragraph{LP construction cost.}
The full LP relaxation contains $|\J||\I|$ shipping variables.
At optimality, the large majority carry positive reduced cost and are
set to zero; they play no role in the LP solution and therefore provide
no useful signal to KS.
Constructing the full LP matrix is therefore unnecessary: column generation
recovers the same LP bound, facility partition, and reduced-cost rankings
without ever adding the irrelevant variables.
On the largest instances (e.g., $|\J|=2{,}000$, $|\I|=2{,}000$), the full LP
solve can dominate total running time, making KS slower than necessary before
the first MILP is even solved.

\paragraph{$\gamma$-threshold arc selection: size and bias.}
G\&S~2012 includes arc $(j,i)$ in a restricted MILP whenever its reduced cost
satisfies $\rc{ji} \leq \gamma$, where $\gamma$ is the mean reduced cost over
LP-active arcs.
Two concerns follow.
First, the number of arcs admitted by $\gamma$ depends on LP concentration and
can vary widely across instances and capacity regimes, giving unpredictable
subproblem sizes.
Second, consider a bucket facility $j \in \Jzero$ (LP dual $\alpha^*_j = 0$).
By equation~\eqref{eq:pricing}, its arc reduced cost is
$\rc{ji} = c_{ji} - \pi^*_i - \alpha^*_j = c_{ji} - \pi^*_i$: the capacity
dual contributes nothing.
For a kernel facility $j' \in \Jone$ with $\alpha^*_{j'} > 0$, the same
formula gives $\rc{j'i} = c_{j'i} - \pi^*_i - \alpha^*_{j'} < c_{j'i} - \pi^*_i$.
The capacity dual reduces kernel arc reduced costs, making them easier to pass
the $\gamma$ threshold.
Bucket arcs must clear a structurally higher hurdle with no capacity signal
subtracted, so fewer bucket arcs qualify for any given $\gamma$.
This treats bucket facilities less favorably than kernel facilities, giving
them a systematically smaller arc set in every restricted MILP.

\paragraph{Cold-start MILP transitions.}
G\&S~2012 passes no information between consecutive restricted MILPs.
When the arc set changes (between stages or between kernel and bucket MILPs),
the incumbent solution from the previous subproblem may not be feasible in the
next one, so the solver must find feasibility from scratch.
In the multi-stage setting this is wasteful: a good solution is available but
cannot be used because there is no guarantee it remains feasible after the
arc set is modified.

FKS addresses each aspect in turn through the components described in the
following sections.

\subsection{Column Generation for the LP}
\label{sec:cg}

\paragraph{Motivation.}
The full LP over all $|\J| \times |\I|$ shipping variables can dominate KS
runtime on large instances.
The MSCFLP LP solution is typically sparse: only a small fraction of shipping
variables carry positive flow at optimality.
Column generation exploits this structure by starting with a small column set and
adding only inactive variables that can improve the LP objective.

\paragraph{Column generation setup.}
We keep all $y_j$ variables in the model from the start (there are only
$|\J|$ of them).
Only the $x_{ji}$ shipping variables are generated on demand.

The \emph{restricted master problem} (RMP) at iteration $t$ contains:
\begin{itemize}
  \item All $y_j$ variables with their opening costs $f_j$.
  \item A subset $\mathcal{A}^t \subseteq \J \times \I$ of active
        shipping columns with costs $c_{ji}$.
  \item Demand constraints~\eqref{eq:demand}: $\sum_{j:(j,i)\in\mathcal{A}^t}
        x_{ji} = d_i$ for each $i$ (initially enforced via a small artificial
        variable to guarantee feasibility).
  \item Capacity constraints~\eqref{eq:capacity}: $\sum_{i:(j,i)\in\mathcal{A}^t}
        x_{ji} \leq q_j y_j$ for each $j$.
  \item VUB constraints~\eqref{eq:vub}: $x_{ji} \leq d_i y_j$ for each
        $(j,i) \in \mathcal{A}^t$.
\end{itemize}
When column $(j,i)$ is added to $\mathcal{A}^t$, its variable $x_{ji}$ enters
the demand constraint for client $i$, the capacity constraint for facility $j$,
and a new VUB constraint.

\paragraph{Initialization.}
We initialize $\mathcal{A}^0$ by selecting the $k_{\text{init}} = 30$
cheapest facilities (by shipping cost $c_{ji}$) for each client $i$.
This gives $30 \times |\I|$ initial columns (120{,}000 for a
$1{,}000 \times 4{,}000$ instance), which is typically sufficient to produce
a feasible RMP.

\paragraph{Pricing step.}
After solving the RMP at iteration $t$, let $\pi_i$ be the dual of demand
constraint $i$ and $\alpha_j$ the dual of capacity constraint $j$.
The reduced cost of any inactive shipping variable $x_{ji}$
($(j,i) \notin \mathcal{A}^t$) is:
\begin{equation}
  \rc{ji} = c_{ji} - \pi_i - \alpha_j.
  \label{eq:pricing}
\end{equation}
If $\rc{ji} < 0$, then adding $x_{ji}$ to the RMP can decrease the LP
objective.
We add \emph{all} inactive columns with $\rc{ji} < 0$ simultaneously, then
re-solve.
This is the ``full pricing'' variant of column generation.
It is more expensive per iteration than pricing a single column, but it avoids
long sequences of small RMP resolves and is effective when many columns have
negative reduced cost in early iterations.
When a shipping column is added, its VUB constraint is introduced at the same
time.
After convergence, reduced costs used for edge selection are taken from the
final LP solution: active generated columns use the solver's reduced costs,
while inactive columns use the final demand and capacity dual prices.
The computational study checks the CG LP against the full LP whenever the full
LP can be solved within the experimental protocol.

\paragraph{Termination.}
When no inactive column has $\rc{ji} < 0$ and all artificial variables are zero,
the current RMP solution is accepted as the LP solution passed to KS.
This follows from LP duality: at the optimal RMP dual $(\pi, \alpha)$,
every column with $\rc{ji} \geq 0$ can only increase the objective if added,
so no ungenerated shipping column is attractive under the pricing oracle.

\paragraph{Vectorized implementation.}
The pricing step~\eqref{eq:pricing} is a matrix subtraction.
Let $C \in \mathbb{R}^{|\J| \times |\I|}$ be the shipping cost matrix.
After solving the RMP, the reduced costs of all $|\J| \times |\I|$ pairs are:
\[
  \mathrm{RC} = C - \boldsymbol{\pi}^{\top}[\text{broadcasted}]
                  - \boldsymbol{\alpha}[\text{broadcasted}].
\]
All inactive entries of $\mathrm{RC}$ with value $< 0$ are added in one
operation.
For $|\J| = 1{,}000$ and $|\I| = 4{,}000$, this is a single NumPy operation on a
$1{,}000 \times 4{,}000$ matrix and is small relative to the RMP solves in our
target setting.

\paragraph{Convergence in practice.}
For each reported instance where the full LP is solved, we compare the CG
objective value, facility partition, and reduced-cost rankings with those
obtained from the full LP.
This check is important because the purpose of CG in FKS is acceleration, not a
change in the information used to guide Kernel Search.
Section~\ref{sec:study_lp} reports the resulting CG iteration counts, active
columns, LP times, and full-LP comparisons.

\paragraph{Remark on branch-and-price.}
Column generation applied to the LP relaxation does \emph{not} make FKS a
branch-and-price algorithm.
Branch-and-price is an exact method that solves an LP at every node of a
branch-and-bound tree~\citep{barnhart1998branch}.
FKS solves the LP exactly once (via CG), then runs heuristic MILPs with no
LP re-solve.
The heuristic character of FKS is unchanged; CG is purely a speed-up for
the single LP solve that KS already requires.

\subsection{Feasibility-Preserving Stage Transitions}
\label{sec:proposition}

When an arc set changes between stages, arcs used by the incumbent may be
absent from the new set, making the incumbent infeasible.
The solver must then rediscover feasibility from scratch.
We eliminate this by including all incumbent arcs in the next stage's edge
set by construction.

Let $(y^s, x^s)$ be the best solution found after Stage~$s$.
For Stage~$s{+}1$ with parameter $k_{s+1}$, we define the edge set as:
\begin{equation}
  \mathcal{E}^{\text{inc}}_{k_{s+1}}(x^s)
  = \bigcup_{i \in \I}
    \Bigl(
      \underbrace{\{(j, i) : x^s_{ji} > 0\}}_{\text{incumbent arcs}}
      \;\cup\;
      \underbrace{\text{top-}k_{s+1}\text{ by } \rc{ji}
      \text{ in } \K}_{\text{reduced-cost fill}}
    \Bigr).
  \label{eq:stage2_edges}
\end{equation}
Each later stage therefore contains the incumbent arcs plus a reduced-cost
fill set.
The edge count is at most $k_{s+1}|\I|$ plus the number of incumbent arcs
not already in the fill set.

\begin{proposition}[Feasibility Preservation]
\label{prop:feasibility}
Let $(y^s, x^s)$ be a feasible solution to the Stage~$s$ restricted MILP
with facility set $\mathcal{F}^s$ and edge set $\mathcal{E}^s$.
Let the Stage~$(s{+}1)$ facility set $\mathcal{F}^{s+1}$ contain every
facility with $y^s_j=1$.
Let $\mathcal{E}^{s+1} = \mathcal{E}^{\mathrm{inc}}_{k_{s+1}}(x^s)$ as in
equation~\eqref{eq:stage2_edges}.
Then $(y^s, x^s)$ is feasible in the Stage~$(s{+}1)$ restricted MILP.
\end{proposition}

\begin{proof}
By construction~\eqref{eq:stage2_edges}, for every client $i$ and every
facility $j$ with $x^s_{ji} > 0$, the arc $(j,i)$ belongs to
$\mathcal{E}^{s+1}$.
The assumption on $\mathcal{F}^{s+1}$ ensures every facility opened by
$y^s$ is available in the next restricted MILP.
We verify the three constraint types:
\begin{enumerate}
  \item \textbf{Demand.} $(y^s, x^s)$ satisfies $\sum_j x^s_{ji} = d_i$
    for every client $i$. The same equality holds in Stage~$(s{+}1)$
    because every arc used by $x^s$ is present.
  \item \textbf{Capacity.} $(y^s, x^s)$ satisfies
    $\sum_i x^s_{ji} \leq q_j y^s_j$. Since $y^s$ and $x^s$ are unchanged,
    this constraint is satisfied.
  \item \textbf{Variable upper bound.} $(y^s, x^s)$ satisfies
    $x^s_{ji} \leq d_i y^s_j$ for every arc with $x^s_{ji} > 0$.
    These arcs are present in $\mathcal{E}^{s+1}$, so the VUB constraints
    are present and satisfied.
\end{enumerate}
Therefore $(y^s, x^s)$ is feasible in Stage~$(s{+}1)$. \hfill$\square$
\end{proof}

\begin{corollary}
\label{cor:warmstart}
When $(y^s, x^s)$ is injected as a MIP start for Stage~$(s{+}1)$, the
solver receives a feasible incumbent by construction and can begin from a
known primal bound.
\end{corollary}

\begin{remark}
Proposition~\ref{prop:feasibility} does not depend on the LP solution
structure, the instance size, or the capacity ratio.
It holds for any multi-stage matheuristic that retains incumbent facilities
and builds edge sets as in equation~\eqref{eq:stage2_edges}, making the
feasibility-preserving construction a general design principle for
Kernel Search algorithms.
\end{remark}

\paragraph{Practical effect.}
Stage~1 uses \texttt{MIPFocus=1} (Gurobi's feasibility-first mode) because
no warm start is yet available.
Stages~2 and~3 inject the previous incumbent and use \texttt{MIPFocus=0}
(balanced), because a primal bound is already known.
This removes the need for later stages to rediscover feasibility after the
arc set has been narrowed.

The feasibility guarantee is also the structural reason the flat-$k$ funnel
in the next section is safe to use: because Proposition~\ref{prop:feasibility}
holds regardless of how narrow the fill set is, the fill size $k$ can be
reduced aggressively across stages without risking cold restarts.

\subsection{Edge Selection: The \texorpdfstring{Flat-$k$}{Flat-k} Criterion}
\label{sec:edge_selection}

Section~\ref{sec:proposition} establishes that all incumbent arcs must be
retained in the next stage's edge set.
This section specifies the remaining arcs: the reduced-cost fill in
equation~\eqref{eq:stage2_edges}.
We use a \emph{client-centric, flat-$k$} criterion: for each client $i$,
include the $k$ facilities with the lowest reduced cost $\rc{ji}$.

\paragraph{Stage~1 edge set.}
Stage~1 has no incumbent, so the edge set is a pure reduced-cost selection:
\begin{equation}
  \mathcal{E}^{\text{RC}}_{k_1}
  = \bigcup_{i \in \I}
    \bigl\{ (j, i) : j \text{ is among the } k_1 \text{ lowest-}\rc{ji}
    \text{ facilities in } \K \bigr\}.
  \label{eq:stage1_edges}
\end{equation}
We use $k_1 = 50$, giving 50{,}000 edges on a $1{,}000 \times 1{,}000$
instance and scaling linearly with $|\I|$ regardless of LP behavior.
Stages~2 and~3 use equation~\eqref{eq:stage2_edges} with $k_2=20$ and
$k_3=10$: the fill set shrinks while the incumbent arcs are always retained
by Proposition~\ref{prop:feasibility}.

\begin{proposition}[Guaranteed subproblem size]
\label{prop:size}
For any MSCFLP instance, any LP solution, and any $k \geq 1$, the
Stage~1 edge set satisfies $|\mathcal{E}^{\mathrm{RC}}_{k_1}| = k_1 \cdot |\I|$.
This bound is exact and holds regardless of LP concentration, capacity
ratio, or instance scale.
\end{proposition}

The proof is immediate: exactly $k_1$ arcs are selected per client and
clients are disjoint.
The Stage~1 subproblem size is fully controlled through $k_1$ alone.

\paragraph{Why flat-$k$ rather than a threshold $\gamma$?}
Section~\ref{sec:audit} gives two reasons.
First, the $\gamma$-threshold provides no guarantee on subproblem size.
Second, it is structurally biased: bucket facilities have $\alpha^*_j = 0$,
so their arc reduced costs receive no capacity-dual discount and face a
higher effective hurdle than kernel-facility arcs under the same $\gamma$.
The flat-$k$ criterion eliminates both problems: it selects exactly $k$ arcs
per client and ranks all facilities by the same reduced-cost ordering,
treating kernel and bucket facilities identically.

\paragraph{Why three stages?}
The specific values $k_1=50$, $k_2=20$, $k_3=10$ are a practical choice,
not a theoretically derived optimum.
What matters is that the sequence is decreasing (so the fill set narrows
across stages) and that the feasibility guarantee of
Proposition~\ref{prop:feasibility} holds at each transition regardless of
how small $k$ becomes.
A practitioner with more time or memory can raise all three values.

\subsection{Full Algorithm}
\label{sec:full_algorithm}

Algorithm~\ref{alg:fks} states FKS precisely.
All parameters not listed in the algorithm follow G\&S~2012 defaults.
We use $k_1 = 50$, $k_2 = 20$, $k_3 = 10$ and a 300-second time limit per
stage throughout.
When a bucket solution improves the incumbent, any opened bucket facility is
promoted into the kernel before the next stage; this keeps the incumbent's
facility set available for the feasibility-preserving transition.
In the pseudocode, $\mathcal{E}^{\text{inc}}_k(x;\mathcal{F})$ denotes the
same incumbent-preserving flat-$k$ construction as
equation~\eqref{eq:stage2_edges}, applied over facility set $\mathcal{F}$.

\begin{algorithm}[t]
\caption{Funnel Kernel Search (FKS)}
\label{alg:fks}
\begin{algorithmic}[1]
\Require Instance $(\J,\I,f,q,d,c)$; stage parameters $(k_s, T_s)_{s=1}^{S}$
\Ensure Feasible integer solution $(y^H, x^H)$

\State \textbf{LP phase:} Solve LP via column generation
       $\to (\xLP, \yLP, \bar{c})$; \quad $z^{\text{LP}} \gets$ LP objective
\State $\K \gets \Jone \cup \Jfrac$
\State Sort $\Jzero$ by $\rc{y_j}$ ascending $\to$ bucket list $L$
\State $z^H \gets +\infty$; \quad $x^{\text{prev}} \gets \emptyset$;
       \quad $y^H \gets \emptyset$

\For{$s = 1, 2, \ldots, S$}

  \State \textbf{Build edge set:}
  \If{$s = 1$}
    \State $\mathcal{E} \gets \mathcal{E}^{\text{RC}}_{k_1}$
           \hfill\Comment{top-$k_1$ by reduced cost per client}
    \State \texttt{MIPFocus} $\gets 1$
           \hfill\Comment{find any feasible incumbent quickly}
  \Else
    \State $\mathcal{E} \gets \mathcal{E}^{\text{inc}}_{k_s}(x^{\text{prev}})$
           \hfill\Comment{preserve incumbent arcs, fill to $k_s$}
    \State \texttt{MIPFocus} $\gets 0$
           \hfill\Comment{balanced; warm start already feasible}
  \EndIf

  \State \textbf{Kernel MILP:}
         $(z, y, x) \gets
         \text{MILP}(\K,\, \mathcal{E},\, z^H,\, \emptyset,\, T_s,\,
         \texttt{warm}=(y^H,x^H))$
  \If{$z < z^H$} \; $z^H, y^H, x^H, x^{\text{prev}} \gets z, y, x, x$
  \EndIf

  \If{$s > 1$} \hfill\Comment{bucket expansion (stages 2+ only)}
    \For{$h = 1, \ldots, \bar{N}_B$}
      \State $B_h \gets$ next $l_{\text{buck}} = |\K|$ facilities from $L$
      \If{$B_h = \emptyset$} \textbf{break} \EndIf
      \State $\mathcal{E}' \gets
             \mathcal{E}^{\text{inc}}_{k_1}(x^H;\,\K \cup B_h)$
             \hfill\Comment{flat-$k$ bucket coverage plus incumbent arcs}
      \State $(z, y, x) \gets
             \text{MILP}(\K \cup B_h,\, \mathcal{E}',\, z^H,\, B_h,\, T_s,\,
             \texttt{warm}=(y^H,x^H))$
      \If{$z < z^H$} \;
        $z^H, y^H, x^H, x^{\text{prev}} \gets z, y, x, x$;
        $\K \gets \K \cup \{j \in B_h : y_j = 1\}$
      \EndIf
    \EndFor
  \EndIf

\EndFor
\State \Return $(y^H, x^H)$
\end{algorithmic}
\end{algorithm}

\paragraph{Why three stages and these $k$ values?}
The specific values $k_1=50$, $k_2=20$, $k_3=10$ are a practical choice, not
a theoretically derived optimum.
The flat-$k$ design is intentionally parameterized: a practitioner with more
time or memory can raise all three values; a practitioner with tighter
constraints can lower them.
What matters structurally is only that the sequence is decreasing (so the
subproblem narrows progressively) and that incumbent arcs are preserved at each
transition (so feasibility is not broken).
The values used here were chosen to give a broad first-stage coverage on the
largest benchmark families ($k=50$ gives 100{,}000 arcs on a $1{,}000\times
2{,}000$ instance) while keeping later stages computationally tractable.
Because every stage after the first begins from a feasible warm start
(Proposition~\ref{prop:feasibility}), later stages spend their time
improving the incumbent rather than rediscovering feasibility.

\paragraph{Relation to G\&S~2012.}
FKS retains all structural components of G\&S~2012: the same kernel
definition, the same bucket expansion structure, and the same MILP
formulation.
The two changes are: (1)~the LP is solved via CG rather than by constructing the
full LP upfront, and (2)~the restricted-MILP phase is replaced by a three-stage
funnel with feasibility-preserving transitions.
The computational study evaluates these changes separately before reporting the
combined end-to-end effect.

\FloatBarrier
\section{Computational Experiments}
\label{sec:experiments}

The computational study uses component-wise comparisons to isolate each of the
three design refinements introduced in Section~\ref{sec:audit}.
In each comparison, exactly one component differs between the two variants,
so the contribution of that specific modification is measured without confounding.
This structure directly mirrors the analysis in Section~\ref{sec:audit}: each
comparison targets one of the three design aspects in isolation.

\subsection{Experimental Setup}

\paragraph{Instances.}
We use the large-scale MSCFLP benchmark families of
\citet{avella2008effective}, also used by \citet{guastaroba2012kernel}.
The main study covers the public Test Bed~A and Test Bed~C families listed in
Table~\ref{tab:instance_families}.
Test Bed~A and Test Bed~C share the same five facility-client scales but use
independently drawn instances, providing a robustness check across instance
variation within each scale.

\begin{table}[htbp]
\centering
\caption{Large-scale benchmark families used in the computational study.}
\label{tab:instance_families}
\small
\setlength{\tabcolsep}{5pt}
\begin{tabular}{llrrrr}
\toprule
Test Bed & Family & $|\J|$ & $|\I|$ & $N$ & Instance indices \\
\midrule
A & p800-4400  &   800 & 4{,}400 & 30 & 1--30 \\
A & p1000-1000 & 1{,}000 & 1{,}000 & 30 & 1--30 \\
A & p1000-4000 & 1{,}000 & 4{,}000 & 30 & 1--30 \\
A & p1200-3000 & 1{,}200 & 3{,}000 & 30 & 1--30 \\
A & p2000-2000 & 2{,}000 & 2{,}000 & 30 & 1--30 \\
\midrule
C & p1000-1000 & 1{,}000 & 1{,}000 & 30 & 61--90 \\
\bottomrule
\end{tabular}
\end{table}

\paragraph{Algorithms compared.}
We report four variants, all implemented in Python~3.11 with Gurobi~13.0.2,
run on an Apple M3~Pro (12-core CPU, 18~GB RAM).
Each instance is solved in an independent process; only one instance runs at a
time, so the full memory and CPU are available to each run.
\begin{itemize}
  \item \textbf{KS-LP}: our re-implementation of G\&S~2012 using the full LP
        relaxation (Gurobi simplex).
        All parameters follow G\&S~2012 defaults.
        KS-LP is the within-study baseline; comparing against it measures
        the combined effect of all three fixes.
  \item \textbf{KS-CG}: the same KS-LP restricted-MILP phase with the full LP
        replaced by column generation.
        Comparing KS-CG against KS-LP isolates Fix~1 (LP cost) alone.
  \item \textbf{KS-CG-WS}: KS-CG with feasibility-preserving stage transitions
        and warm starting added, but retaining the original $\gamma$-threshold
        arc selection.
        Comparing KS-CG-WS against KS-CG isolates Fix~3 alone.
  \item \textbf{FKS-CG-WS}: Funnel Kernel Search --- KS-CG-WS with the
        $\gamma$-threshold replaced by the flat-$k$ criterion,
        $(k_1, k_2, k_3) = (50, 20, 10)$.
        Comparing FKS-CG-WS against KS-CG-WS isolates Fix~2 alone.
        Comparing FKS-CG-WS against KS-LP gives the combined end-to-end picture.
\end{itemize}

\paragraph{Implementation details.}
KS-LP uses the parameter choices of G\&S~2012: the LP-defined initial kernel,
bucket size $l_{\text{buck}} = |\K|$, at most three analyzed buckets ($\bar{N}_B = 3$),
objective cutoffs, and the bucket-forcing constraint.
FKS uses the same kernel and bucket logic but replaces the single restricted-MILP
phase with the flat-$k$ funnel described in Section~\ref{sec:algorithm}.
The restricted-MILP time limit is 300~seconds per subproblem for all variants.
G\&S~2012 allocated 900~seconds per subproblem~\citep{guastaroba2012kernel}
(Section~3.2), noting that ``the software often finds quickly an optimal
or near-optimal solution'' and that the limit exists mainly as a safeguard.
Modern hardware and Gurobi reduce per-subproblem solve times considerably
relative to the CPLEX~12.2 environment of G\&S~2012, so 300~seconds proves
sufficient: our KS-LP recreation matches or slightly exceeds their published
aggregate results on p1000-1000 (Section~\ref{sec:literature}).
Because all four variants share this 300-second limit, all within-study
comparisons are unaffected by the difference from the original paper.
Runs are recorded with instance name, test bed, LP solver, stage widths,
LP time, restricted-MILP time, objective value, and gap.

\paragraph{Metrics.}
The LP gap is
\[
  \text{gap}(\%) = \frac{z^H - z^{\mathrm{LP}}}{z^{\mathrm{LP}}} \times 100,
\]
where $z^H$ is the heuristic objective and $z^{\mathrm{LP}}$ is the LP lower
bound from the same run.
Because our LP lower bounds can differ slightly from those reported by
G\&S~2012 due to solver and hardware differences (see Section~\ref{sec:lp_bounds}),
we also report integer objectives for direct literature comparisons.
Win/tie/loss counts use a tolerance of 0.01 cost units.
Runtime is decomposed into LP time, restricted-MILP time, and total time
whenever the breakdown is relevant.

\subsection{Study 1: LP Bottleneck Removal}
\label{sec:study_lp}

The first study measures how much of the KS setup time is consumed by the LP
phase and how accurately column generation replicates the full LP solution.
This is primarily a timing and fidelity check, not a heuristic-quality
comparison.

Table~\ref{tab:lp_bottleneck} summarizes LP timing across all benchmark families.
Two patterns stand out.
First, CG LP always matches the full LP objective to machine precision: the
two lower bounds agree in all reported significant figures on every completed
instance, confirming that column generation supplies the same information to KS
without altering its guidance.
Second, the CG speedup grows sharply with instance size.
For p1000-1000 (Test Bed~A), the full LP averages 58.0~seconds and CG reduces
this to 1.3~seconds (45.6$\times$) using 3.4~iterations and activating only 3.3\%
of the $10^6$ possible shipping arcs.
For p2000-2000, the full LP averages 710.7~seconds and CG reduces this to
4.2~seconds, a 170.2$\times$ speedup.
The active-arc percentage falls as the instance grows (3.8\% for p800-4400 to
2.4\% for p2000-2000), consistent with larger instances having sparser LP
solutions and therefore more redundant columns to avoid.

\begin{table}[htbp]
\centering
\caption{LP relaxation: full LP (Gurobi simplex) versus column-generation LP.
         Speedup is the ratio of average full-LP time to average CG-LP time.
         Active arcs reports the percentage of all $|\J|\times|\I|$ arcs
         variables present in the terminal CG LP.
         CG LP objectives match full LP objectives to machine precision on all
         completed instances.}
\label{tab:lp_bottleneck}
\footnotesize
\setlength{\tabcolsep}{3pt}
\begin{tabular}{llrrrrr}
\toprule
Test Bed & Family & Full LP (s) & CG LP (s) & Speedup & CG iters & Active arcs (\%) \\
\midrule
A & p800-4400  & 276.5 & 5.9 & 46.7$\times$  & 2.6 & 3.8 \\
A & p1000-1000 &  58.0 & 1.3 & 45.6$\times$  & 3.4 & 3.3 \\
A & p1000-4000 & 409.1 & 6.9 & 59.3$\times$  & 2.8 & 3.0 \\
A & p1200-3000 & 389.3 & 4.5 & 85.8$\times$  & 3.3 & 2.5 \\
A & p2000-2000 & 710.7 & 4.2 & 170.2$\times$ & 3.9 & 2.4 \\
\midrule
C & p1000-1000 &   8.8 & 0.7 & 13.5$\times$  & 1.3 & 3.0 \\
\bottomrule
\end{tabular}
\end{table}

\subsubsection{On LP Lower-Bound Differences Between Studies}
\label{sec:lp_bounds}

Our full LP lower bounds for p1000-1000 match those of G\&S~2012 exactly on
tight-capacity instances (instances~1--5), but are slightly higher on
looser-capacity instances (instances~16--30), by up to 25~cost units.
We treat these differences as numerical and reporting differences across
software generations rather than as an algorithmic improvement in the LP bound.
The original study used a different solver environment and reports rounded
values, while our within-study comparisons use LP bounds computed in the same
Gurobi environment for all three variants.
The important validation for this paper is therefore internal: the CG LP
reproduces our full LP lower bound to within $1.3 \times 10^{-5}$ on all
instances, confirming that column generation supplies the same LP information
used by the full-LP KS baseline.

\subsection{Study 2: Restricted-MILP Acceleration}
\label{sec:study_milp}

The second study compares KS-CG and FKS-CG-WS on the restricted-MILP phase.
Both variants receive identical CG LP output, so any difference in MILP time
or solution quality is attributable to the flat-$k$ arc selection and
feasibility-preserving stage transitions combined.
Table~\ref{tab:kscg_fkscg} reports results across all families; a separate
component-wise breakdown isolating each fix individually is given in
Section~\ref{sec:study_components}.

FKS-CG-WS never loses on solution quality: across all five Test Bed~A families
(150 instances total), FKS-CG-WS matches or improves KS-CG on every instance.
The quality advantage is most pronounced on p1000-4000 (23W/7T/0L) and
p2000-2000 (22W/8T/0L), where the flat-$k$ arc selection appears to
consistently deliver better facility combinations than the $\gamma$-threshold.

MILP-phase speedup varies by family and depends on how incumbents interact with
stage transitions.
On families where LP solutions are relatively concentrated (p1000-1000: 1.86$\times$;
p2000-2000: 2.18$\times$), the funnel's narrowing edge sets reduce subproblem size
at each stage and FKS-CG-WS is substantially faster.
On families with more diffuse LP support (p800-4400: 0.94$\times$; p1000-4000:
0.80$\times$), incumbent-arc preservation fills later-stage edge sets with arcs
from the Stage~1 solution, limiting the thinning effect of the funnel.
In these cases FKS-CG-WS's MILP phase is slightly slower than KS-CG, but end-to-end
FKS-CG-WS remains faster once the LP-phase savings are included (Table~\ref{tab:end_to_end}).

Test Bed~C p1000-1000 (30 independent instances, indices 61--90) shows
1.76$\times$ MILP speedup with 7W/22T/1L.
The single loss is instance~82, where KS-CG finds an objective 2.6 cost units
lower (a 0.00007\% relative difference), consistent with numerical noise rather
than a systematic quality difference.

\begin{table}[htbp]
\centering
\caption{Restricted-MILP phase: KS-CG versus FKS-CG-WS on the same CG LP.
         MILP time excludes the shared LP phase.
         Speedup is KS-CG time divided by FKS-CG-WS time.
         W/T/L counts wins, ties, and losses for FKS-CG-WS relative to KS-CG
         by integer objective ($\pm 0.01$ tolerance).}
\label{tab:kscg_fkscg}
\footnotesize
\setlength{\tabcolsep}{3.5pt}
\begin{tabular}{llrrrrr}
\toprule
 & & & \multicolumn{2}{c}{Avg.\ MILP time (s)} & & \\
\cmidrule(lr){4-5}
TB & Family & $N$ & KS-CG & FKS-CG-WS & Speedup & W/T/L \\
\midrule
A & p800-4400  & 30 & 513.7 & 546.3 & 0.94$\times$ & 12/18/0 \\
A & p1000-1000 & 30 &  80.6 &  43.4 & 1.86$\times$ &  1/29/0 \\
A & p1000-4000 & 30 & 444.3 & 555.0 & 0.80$\times$ & 23/7/0  \\
A & p1200-3000 & 30 & 508.8 & 481.3 & 1.06$\times$ & 16/14/0 \\
A & p2000-2000 & 30 & 479.2 & 219.5 & 2.18$\times$ & 22/8/0  \\
\midrule
C & p1000-1000 & 30 & 125.1 &  71.2 & 1.76$\times$ &  7/22/1 \\
\bottomrule
\end{tabular}
\end{table}

\subsection{Study 3: End-to-End Performance}
\label{sec:study_end_to_end}

Table~\ref{tab:end_to_end} reports end-to-end results across all families.
FKS-CG-WS never loses on solution quality against KS-LP across all 150 Test Bed~A
instances: the W/T/L counts show wins or ties only.
The same holds against KS-CG except for the single numerically negligible
loss on Test Bed~C instance~82 noted above.

The end-to-end speedup over KS-LP grows with instance size, driven primarily
by the LP phase.
For p1000-1000, where the full LP is already modest (58~s), the combined speedup
is 3.08$\times$.
For p2000-2000, where the full LP averages 710.7~s and column generation
reduces this to 4.2~s, the combined speedup reaches 5.32$\times$.

The pattern across families also illustrates the interaction between the two
phases.
On p800-4400, FKS-CG-WS's MILP phase is slightly slower than KS-CG (0.94$\times$),
but the LP-phase saving of 270~s per instance still produces a 1.52$\times$
end-to-end gain over KS-LP.
On p1000-4000, FKS-CG-WS's MILP phase is again slower than KS-CG (0.80$\times$),
yet end-to-end FKS-CG-WS is 2.01$\times$ faster than KS-LP because the LP saving
(409~s full LP versus 6.9~s CG) dominates.
On p2000-2000, both phases favor FKS-CG-WS simultaneously: the LP phase contributes
a 170$\times$ speedup and the MILP phase contributes a further 2.18$\times$,
yielding the largest combined gain in the study.

\begin{table}[htbp]
\centering
\caption{End-to-end total time (LP + restricted-MILP) for all three variants.
         Speedup is KS-LP time divided by FKS-CG-WS time.
         W/T/L counts wins, ties, and losses for FKS-CG-WS relative to each
         baseline by integer objective ($\pm 0.01$ tolerance).
         Solution quality (gap) is reported separately in
         Table~\ref{tab:kscg_fkscg}.}
\label{tab:end_to_end}
\footnotesize
\setlength{\tabcolsep}{3pt}
\begin{tabular}{llrrrrrrr}
\toprule
 & & & \multicolumn{3}{c}{Avg.\ total time (s)} & & \multicolumn{2}{c}{W/T/L (FKS-CG-WS vs)} \\
\cmidrule(lr){4-6}\cmidrule(lr){8-9}
TB & Family & $N$ & KS-LP & KS-CG & FKS-CG-WS & Speedup & KS-LP & KS-CG \\
\midrule
A & p800-4400  & 30 &   838.3 & 519.6 & 552.3 & 1.52$\times$ & 15/15/0 & 12/18/0 \\
A & p1000-1000 & 30 &   137.9 &  81.8 &  44.7 & 3.08$\times$ &  4/26/0 &  1/29/0 \\
A & p1000-4000 & 30 & 1{,}131.7 & 449.9 & 561.9 & 2.01$\times$ & 17/13/0 & 23/7/0  \\
A & p1200-3000 & 30 & 1{,}048.6 & 513.3 & 485.8 & 2.16$\times$ & 16/14/0 & 16/14/0 \\
A & p2000-2000 & 30 & 1{,}189.3 & 483.6 & 223.7 & 5.32$\times$ & 22/8/0  & 22/8/0  \\
\midrule
C & p1000-1000 & 30 & 185.1 & 125.8 & 71.8 & 2.58$\times$ & 4/25/1 & 7/22/1 \\
\bottomrule
\end{tabular}
\end{table}

\subsection{Study 4: Component-wise Isolation of Fix~2 and Fix~3}
\label{sec:study_components}

Studies 1--3 compare FKS-CG-WS against KS-CG, which differs in two ways: it uses
flat-$k$ arc selection (Fix~2) and feasibility-preserving stage transitions
(Fix~3).
To isolate each fix, we compare against KS-CG-WS, which adds only the
feasibility-preserving transitions to KS-CG while retaining the $\gamma$-threshold.

\begin{itemize}
  \item KS-CG versus KS-CG-WS: measures Fix~3 (feasibility-preserving
        transitions) alone, with arc selection held constant.
  \item KS-CG-WS versus FKS-CG-WS: measures Fix~2 (flat-$k$ arc selection) alone,
        with feasibility-preserving transitions held constant.
\end{itemize}

Results for p1000-4000 Test Bed~A (30 instances) will be reported for this
ablation because its size and capacity regime produce clear quality differences
in Study~2.
The comparison uses the same instance set, CG LP output, and 300-second
restricted-MILP limit as the main study.
We keep this ablation separate from Tables~\ref{tab:kscg_fkscg}
and~\ref{tab:end_to_end} because it requires the additional KS-CG-WS control
run; those runs are not mixed into the aggregate tables until the complete
family is available.

\subsection{Study 5: Comparison Against Proven Optima}
\label{sec:study_proven_optima}

The final study evaluates solution quality against the best available results
from the literature for exact methods.
For Test Bed~A, we use the proven optimal values reported by
\citet{caserta2020corridor}, who solved 110 of the 150 instances to optimality.
For Test Bed~C, we use the proven optimal values and best-known upper bounds
from \citet{avella2021weakflow}, who provide the current state-of-the-art for
these harder instances.
The gap is computed as $\text{gap}(\%) = 100 \times (z^H - z^*) / z^*$, where
$z^*$ is the proven optimal or best-known objective.

Table~\ref{tab:proven_optima_gap} reports the average and maximum gaps for the
110 Test Bed~A instances with known optima.
The results show that FKS-CG-WS is exceptionally effective, finding the proven
optimal solution on 109 of 110 instances.
The single non-optimal solution (for instance p2000-2000-17) has a gap of just
0.0007\%, corresponding to an objective value less than one cost unit from the
optimum.
This confirms that the FKS design not only accelerates the search but also
consistently guides it to solutions of very high quality.

In contrast, the baseline KS-LP and KS-CG methods, while still strong, show
small but consistent gaps across all families.
Notably, KS-CG produces a large outlier gap of 0.27\% on one instance
(p1000-4000-17), where FKS-CG-WS finds the optimum. This suggests the flat-$k$
arc selection criterion is more robust than the $\gamma$-threshold.

\begin{table}[htbp]
\centering
\caption{Gap (\%) versus proven optimal solutions on 110 Test Bed~A instances.
         FKS-CG-WS finds the proven optimum on 109 of 110 instances.}
\label{tab:proven_optima_gap}
\footnotesize
\setlength{\tabcolsep}{4pt}
\begin{tabular}{lrrrrrr}
\toprule
& \multicolumn{2}{c}{KS-LP} & \multicolumn{2}{c}{KS-CG} & \multicolumn{2}{c}{FKS-CG-WS} \\
\cmidrule(lr){2-3}\cmidrule(lr){4-5}\cmidrule(lr){6-7}
Family (N) & Avg.\ gap & Max gap & Avg.\ gap & Max gap & Avg.\ gap & Max gap \\
\midrule
p800-4400 (21)   & 0.0021\% & 0.0246\% & 0.0010\% & 0.0055\% & 0.0000\% & 0.0000\% \\
p1000-1000 (30)  & 0.0018\% & 0.0373\% & 0.0001\% & 0.0025\% & 0.0000\% & 0.0003\% \\
p1000-4000 (21)  & 0.0029\% & 0.0173\% & 0.0211\% & 0.2726\% & 0.0000\% & 0.0000\% \\
p1200-3000 (19)  & 0.0027\% & 0.0268\% & 0.0014\% & 0.0096\% & 0.0000\% & 0.0000\% \\
p2000-2000 (19)  & 0.0026\% & 0.0136\% & 0.0012\% & 0.0088\% & 0.0000\% & 0.0007\% \\
\midrule
\textbf{Overall (110)} & \textbf{0.0024\%} & \textbf{0.0373\%} & \textbf{0.0047\%} & \textbf{0.2726\%} & \textbf{0.0000\%} & \textbf{0.0007\%} \\
\bottomrule
\end{tabular}
\end{table}

\subsection{Validating the KS-LP Baseline Against Published Results}
\label{sec:literature}

\paragraph{Why a re-implementation is necessary.}
The computational environments in G\&S~2012 and this study differ in solver
(CPLEX~12.2 versus Gurobi), hardware (Intel Xeon 2.27~GHz versus modern
hardware), and per-subproblem time limit (900~seconds versus 300~seconds).
These differences make direct comparisons of running times meaningless, and
they mean we cannot use the published G\&S timing or solution results as the
control in our ablation study.
Instead, we implement a faithful G\&S-style baseline --- same formulation, same
parameter choices ($\gamma$, $l_{\text{buck}}$, $\bar{N}_B$, $p$), same
kernel construction and bucket logic --- and run it as KS-LP in our own
environment.
All proposed variants (KS-CG, KS-CG-WS, FKS-CG-WS) are then compared against
this same KS-LP baseline, ensuring that solver, hardware, and time limit are
identical across every comparison in the study.

\paragraph{How we know the re-implementation is fair.}
The validity of this design depends entirely on whether our KS-LP recreation
is a faithful and adequate representation of the original algorithm.
If our re-implementation were systematically weaker than G\&S's, the study
would be unfairly easy: our new methods would beat a handicapped baseline.
We therefore validate KS-LP against G\&S's published results on the p1000-1000
Test Bed~A family, using the same benchmark quantity G\&S report: improvement
over the Avella et al.~\citep{avella2008effective} published upper bounds.

G\&S report an average improvement of $-0.13\%$ over those bounds, with 26
improvements out of 30 instances.
Our KS-LP achieves the same average improvement figure, improves 28 upper
bounds, and ties the remaining 2 --- despite using only one-third of
their per-subproblem time budget.
Because G\&S report rounded optimality gaps rather than raw integer objectives,
we also perform a secondary reconstruction check using
\[
  \tilde{z}^{H}_{\mathrm{BKS}}
  =
  \frac{z^{LB}_{\mathrm{BKS}}}
       {1-\mathrm{Gap}_{\mathrm{BKS}}/100},
\]
which gives the same aggregate conclusion: our KS-LP objective is slightly
lower than the reconstructed B-KS average (treated as an approximate
consistency check only, given the rounded inputs).

Our KS-LP re-implementation therefore produces \emph{at least as good}
solution quality as the published G\&S results on this family.
This confirms that KS-LP is a fair, non-handicapped representation of
the original algorithm, and that any gains attributed to KS-CG, KS-CG-WS,
or FKS-CG-WS in the following studies represent genuine improvements over
a credible baseline.
Detailed per-instance results appear in Appendix~\ref{app:detail}.

\FloatBarrier
\section{Conclusion}
\label{sec:conclusion}

Kernel Search for the large-scale MSCFLP is effective because it uses the LP
relaxation to identify a small set of promising facilities and arcs before
solving restricted MILPs.
The same dependence on the LP, however, exposes a bottleneck on the largest
instances: obtaining the LP information can itself become expensive.
This paper addresses that bottleneck directly by solving the KS LP relaxation
through column generation.
The intention is not to alter the heuristic guidance, but to obtain the
LP-derived bound, facility partition, and reduced-cost signal at lower setup
cost.

The second contribution concerns the restricted-MILP phase.
The original KS framework already contains several speed controls, such as
limiting bucket exploration and removing unpromising facilities from the kernel.
Those controls reduce MILP effort but naturally introduce a speed-quality
trade-off.
FKS instead organizes the restricted search as a flat-$k$ funnel.
The flat-$k$ rule gives predictable base subproblem sizes, while preserving all
incumbent arcs and facilities between consecutive stages guarantees that the
previous solution is feasible in the next restricted MILP.
This gives a simple structural basis for warm-starting later stages after the
edge set has been narrowed.

The broader message is diagnostic.
Matheuristic improvements benefit from examining the existing method's design
choices carefully before proposing new components.
For G\&S~2012 KS on the MSCFLP, a close reading surfaces three design aspects
that can each be refined --- avoidable LP cost, biased arc selection, and
cold-start MILP transitions --- each with a targeted, independently verifiable
modification.
The computational study is organized around this structure: each comparison
changes exactly one component, making it possible to attribute improvements
to a specific design decision rather than to the aggregate effect of the method.

Several extensions remain natural.
First, completing the remaining Test Bed~A and Test Bed~C families will provide
a broader robustness check across facility-client scales and capacity regimes.
Second, per-client adaptive budgets could allocate more arcs to clients with
diffuse LP support and fewer arcs to clients with concentrated support while
retaining the incumbent-preservation guarantee.
Third, the feasibility-preservation principle may be useful in other Kernel
Search settings whenever restricted subproblems are built by selecting subsets
of variables used by an incumbent solution.


\bibliographystyle{elsarticle-harv}
\bibliography{references}

\appendix

\section{Detailed Per-Instance Results}
\label{app:detail}

Tables~\ref{tab:app_tba_1000_1000} and~\ref{tab:app_tba_800_4400} give
per-instance results for the p1000-1000 and p800-4400 Test Bed~A families
respectively.
Table~\ref{tab:app_tba_1000_1000} gives per-instance results for the
p1000-1000 Test Bed~A family.
The LP lower bound $z^{\mathrm{LP}}$ is shared by all three variants (the CG LP
reproduces the full LP objective to machine precision on every instance).
Objective values and $z^{\mathrm{LP}}$ allow the reader to verify the reported
gap directly as $(z^H - z^{\mathrm{LP}})/z^{\mathrm{LP}} \times 100$.
Time is total wall-clock seconds (LP + restricted-MILP).
A bold FKS-CG objective indicates a strictly better solution than KS-LP.

\begin{table}[htbp]
\centering
\caption{Per-instance results: p1000-1000, Test Bed~A.
         $z^{\mathrm{LP}}$: LP lower bound (identical for all three variants).
         Objective columns report integer objective values;
         time columns report total wall-clock seconds (LP + MILP).
         Bold FKS-CG objective indicates a strictly better solution than KS-LP.}
\label{tab:app_tba_1000_1000}
\footnotesize
\setlength{\tabcolsep}{3pt}
\begin{tabular}{lrrrrrrr}
\toprule
 & & \multicolumn{2}{c}{KS-LP} & \multicolumn{2}{c}{KS-CG} & \multicolumn{2}{c}{FKS-CG} \\
\cmidrule(lr){3-4}\cmidrule(lr){5-6}\cmidrule(lr){7-8}
Instance & $z^{\mathrm{LP}}$ & Obj. & Time & Obj. & Time & Obj. & Time \\
\midrule
p1000-1000-1  &  920{,}490.9 &  920{,}505.2 & 117.0 &  920{,}505.2 &  64.6 &  920{,}505.2 &  42.4 \\
p1000-1000-2  &  920{,}809.1 &  920{,}834.2 &  88.4 &  920{,}834.2 &  37.2 &  920{,}834.2 &  10.5 \\
p1000-1000-3  &  931{,}266.1 &  931{,}305.4 & 109.5 &  931{,}305.4 &  69.1 &  931{,}305.4 &  17.3 \\
p1000-1000-4  &  917{,}924.5 &  917{,}951.7 & 114.7 &  917{,}951.7 &  65.9 &  917{,}951.7 &  10.2 \\
p1000-1000-5  &  915{,}609.1 &  915{,}635.4 & 118.1 &  915{,}635.4 &  64.7 &  915{,}635.4 &  12.2 \\
p1000-1000-6  &  530{,}717.7 &  530{,}737.5 & 245.9 &  530{,}737.5 & 138.6 &  530{,}737.5 &  30.9 \\
p1000-1000-7  &  570{,}761.3 &  570{,}792.9 & 340.0 &  570{,}792.9 & 292.2 &  570{,}792.9 &  40.7 \\
p1000-1000-8  &  573{,}812.7 &  573{,}827.8 & 131.3 &  573{,}827.8 &  61.3 &  573{,}827.8 &  10.7 \\
p1000-1000-9  &  556{,}659.9 &  556{,}686.4 & 167.6 &  556{,}686.4 &  91.3 &  556{,}686.4 &  14.0 \\
p1000-1000-10 &  543{,}588.6 &  543{,}599.8 & 149.4 &  543{,}599.8 &  57.5 &  543{,}599.8 &   7.2 \\
p1000-1000-11 &  350{,}058.3 &  350{,}069.3 & 115.6 &  350{,}069.3 &  76.5 &  350{,}069.3 &  13.2 \\
p1000-1000-12 &  311{,}493.4 &  311{,}516.2 & 131.3 &  311{,}516.2 &  52.7 &  311{,}516.2 &  12.5 \\
p1000-1000-13 &  331{,}134.3 &  331{,}149.2 & 119.3 &  331{,}149.2 &  48.2 &  331{,}149.2 &  11.1 \\
p1000-1000-14 &  334{,}602.1 &  334{,}629.8 & 164.3 &  334{,}630.9 &  74.6 & \textbf{334{,}623.6} &  14.9 \\
p1000-1000-15 &  347{,}379.5 &  347{,}407.3 & 144.6 &  347{,}407.3 &  72.8 &  347{,}407.3 &  15.1 \\
p1000-1000-16 &   95{,}940.6 &   95{,}959.2 & 113.2 &   95{,}959.2 &  37.8 &   95{,}959.2 &  15.4 \\
p1000-1000-17 &   93{,}969.0 &   93{,}998.8 &  92.6 &   93{,}998.8 &  38.1 &   93{,}998.8 &  13.6 \\
p1000-1000-18 &   95{,}711.8 &   95{,}741.1 & 115.9 &   95{,}737.2 &  77.6 &  \textbf{95{,}737.2} &  30.4 \\
p1000-1000-19 &   91{,}028.2 &   91{,}057.3 & 128.9 &   91{,}057.3 &  58.2 &   91{,}057.3 &  22.2 \\
p1000-1000-20 &   97{,}265.2 &   97{,}280.1 &  79.7 &   97{,}280.1 &  34.7 &   97{,}280.1 &  15.5 \\
p1000-1000-21 &   49{,}845.2 &   49{,}884.8 &  89.7 &   49{,}884.8 &  23.6 &   49{,}884.8 &  16.6 \\
p1000-1000-22 &   51{,}086.2 &   51{,}137.5 & 147.4 &   51{,}137.5 &  95.1 &   51{,}137.5 &  46.3 \\
p1000-1000-23 &   50{,}015.9 &   50{,}066.5 & 118.7 &   50{,}060.8 &  43.9 &  \textbf{50{,}060.8} &  30.3 \\
p1000-1000-24 &   50{,}437.8 &   50{,}504.3 & 178.8 &   50{,}504.3 &  72.8 &   50{,}504.3 &  55.8 \\
p1000-1000-25 &   48{,}247.2 &   48{,}338.1 & 282.4 &   48{,}338.1 & 343.2 &   48{,}338.1 & 283.1 \\
p1000-1000-26 &   27{,}364.8 &   27{,}491.0 & 117.3 &   27{,}491.0 &  78.0 &   27{,}491.0 & 136.7 \\
p1000-1000-27 &   26{,}518.5 &   26{,}586.9 & 118.4 &   26{,}586.9 &  79.9 &   26{,}586.9 &  94.3 \\
p1000-1000-28 &   26{,}471.6 &   26{,}563.9 &  98.2 &   26{,}563.9 &  72.7 &   26{,}563.9 &  61.6 \\
p1000-1000-29 &   27{,}758.3 &   27{,}869.8 & 112.0 &   27{,}859.4 &  80.8 &  \textbf{27{,}859.4} & 182.7 \\
p1000-1000-30 &   26{,}740.6 &   26{,}825.2 &  85.5 &   26{,}825.2 &  50.0 &   26{,}825.2 &  73.8 \\
\midrule
Average       &              &              & 137.9 &              &  81.8 &              &  44.7 \\
\bottomrule
\end{tabular}
\end{table}

\begin{table}[htbp]
\centering
\caption{Per-instance results: p800-4400, Test Bed~A.
         $z^{\mathrm{LP}}$: LP lower bound (identical for all three variants).
         Objective columns report integer objective values;
         time columns report total wall-clock seconds (LP + MILP).
         Bold FKS-CG objective indicates a strictly better solution than KS-LP.
         Instances~1--15 are tight-capacity ($r \leq 2$);
         instances~16--30 are loose-capacity ($r \geq 3$).}
\label{tab:app_tba_800_4400}
\footnotesize
\setlength{\tabcolsep}{3pt}
\begin{tabular}{lrrrrrrr}
\toprule
 & & \multicolumn{2}{c}{KS-LP} & \multicolumn{2}{c}{KS-CG} & \multicolumn{2}{c}{FKS-CG} \\
\cmidrule(lr){3-4}\cmidrule(lr){5-6}\cmidrule(lr){7-8}
Instance & $z^{\mathrm{LP}}$ & Obj. & Time & Obj. & Time & Obj. & Time \\
\midrule
p800-4400-1  & 745{,}124.2 & 745{,}157.6 &   739.8 & 745{,}157.6 &   146.0 & 745{,}157.6 &    90.4 \\
p800-4400-2  & 742{,}319.9 & 742{,}345.9 &   690.7 & 742{,}345.9 &   474.4 & 742{,}345.9 &    98.5 \\
p800-4400-3  & 747{,}233.2 & 747{,}284.7 &   601.1 & 747{,}295.7 &   366.1 & 747{,}284.7 &    74.3 \\
p800-4400-4  & 754{,}867.3 & 754{,}884.2 &   479.5 & 754{,}884.2 &   272.6 & 754{,}884.2 &    69.2 \\
p800-4400-5  & 749{,}451.5 & 749{,}498.3 &   769.0 & 749{,}490.7 &   549.1 & \textbf{749{,}490.7} &    79.6 \\
p800-4400-6  & 432{,}514.1 & 432{,}597.8 &   758.4 & 432{,}597.8 &   530.0 & \textbf{432{,}588.2} &   288.5 \\
p800-4400-7  & 447{,}371.0 & 447{,}436.8 & 1149.8 & 447{,}436.8 &   555.3 & 447{,}436.8 &   254.2 \\
p800-4400-8  & 468{,}666.4 & 468{,}724.7 &   837.7 & 468{,}724.7 &   554.1 & 468{,}724.7 &   130.7 \\
p800-4400-9  & 424{,}817.2 & 424{,}888.3 &   803.8 & 424{,}888.3 &   514.9 & \textbf{424{,}876.0} &   116.7 \\
p800-4400-10 & 441{,}145.9 & 441{,}234.7 &   783.3 & 441{,}234.7 &   472.2 & \textbf{441{,}218.5} &   135.3 \\
p800-4400-11 & 290{,}454.5 & 290{,}562.1 &   832.0 & 290{,}562.1 &   461.0 & \textbf{290{,}546.2} &   139.8 \\
p800-4400-12 & 292{,}151.3 & 292{,}220.8 &   733.5 & 292{,}220.8 &   421.9 & \textbf{292{,}210.1} &   194.1 \\
p800-4400-13 & 304{,}467.3 & 304{,}561.7 &   828.3 & 304{,}561.7 &   498.4 & 304{,}561.7 &   275.4 \\
p800-4400-14 & 304{,}332.3 & 304{,}417.4 &   827.6 & 304{,}417.4 &   458.6 & 304{,}417.4 &   203.4 \\
p800-4400-15 & 282{,}818.8 & 282{,}870.4 &   820.4 & 282{,}870.4 &   341.3 & 282{,}870.4 &    78.9 \\
\midrule
p800-4400-16 & 100{,}510.0 & 100{,}662.7 &   663.8 & 100{,}662.7 &   419.9 & 100{,}662.7 &   704.9 \\
p800-4400-17 &  99{,}790.5 &  99{,}985.7 & 1898.2 & 100{,}018.3 &   667.5 & \textbf{99{,}966.9} & 1161.5 \\
p800-4400-18 & 103{,}716.6 & 103{,}945.4 & 1244.9 & 103{,}962.2 &   703.4 & \textbf{103{,}921.1} & 1016.1 \\
p800-4400-19 & 104{,}544.4 & 104{,}689.0 &   783.8 & 104{,}689.0 &   468.5 & 104{,}689.0 &   743.7 \\
p800-4400-20 & 102{,}307.9 & 102{,}456.4 &   576.2 & 102{,}453.9 &   485.0 & \textbf{102{,}453.9} &   374.0 \\
p800-4400-21 &  72{,}364.3 &  72{,}593.0 &   842.5 &  72{,}558.3 &   841.3 & \textbf{72{,}539.1} & 1236.7 \\
p800-4400-22 &  70{,}552.4 &  70{,}829.3 & 1413.2 &  70{,}823.9 &   910.0 & \textbf{70{,}823.9} & 1432.1 \\
p800-4400-23 &  72{,}483.0 &  72{,}796.8 & 1540.4 &  72{,}773.3 &   578.9 & \textbf{72{,}773.3} & 1712.2 \\
p800-4400-24 &  72{,}665.4 &  73{,}032.5 & 1423.5 &  73{,}018.0 &   999.5 & \textbf{73{,}005.6} & 1611.8 \\
p800-4400-25 &  74{,}205.7 &  75{,}091.9 &   416.3 &  74{,}540.5 & 1138.9 & \textbf{74{,}537.5} & 1530.6 \\
p800-4400-26 &  59{,}451.0 &  59{,}499.3 &   227.6 &  59{,}499.3 &   269.1 &  59{,}499.3 &   453.3 \\
p800-4400-27 &  58{,}738.3 &  58{,}866.6 &   394.8 &  58{,}852.1 &   503.9 & \textbf{58{,}852.1} &   481.0 \\
p800-4400-28 &  59{,}954.1 &  60{,}063.0 & 1327.7 &  60{,}063.0 &   559.7 &  60{,}063.0 &   962.4 \\
p800-4400-29 &  57{,}817.3 &  57{,}936.4 &   491.0 &  57{,}938.5 &   276.9 &  57{,}936.4 &   692.3 \\
p800-4400-30 &  57{,}973.0 &  58{,}041.4 &   251.0 &  58{,}041.4 &   150.0 &  58{,}041.4 &   226.0 \\
\midrule
Average       &             &             &   838.3 &             &   519.6 &             &   552.3 \\
\bottomrule
\end{tabular}
\end{table}


\end{document}
